\documentclass[10pt,a4paper]{amsart}
\usepackage[margin=19mm]{geometry}
\usepackage{amsmath,amssymb,mathtools}
\usepackage[scr=rsfs,cal=euler]{mathalfa}
\usepackage{xfrac}
\usepackage[unicode,hidelinks,bookmarks=false]{hyperref}
\newcommand{\ie}{i.e.\ }
\newcommand{\cf}{cf.\ }
\newcommand{\resp}{respectively\ }
\newcommand{\Subsection}{\subsection}
\providecommand{\nobreakdash}{\nobreak\hbox{-}}
\setlength{\emergencystretch}{2em}
\multlinegap0pt
\usepackage{bm}
\usepackage{bbm}
\usepackage{dsfont}
\usepackage{xspace}
\usepackage{cancel}
\usepackage{mathtools}
\usepackage{amsthm}
\makeatletter
\@ifundefined{lemma}{\newtheorem{lemma}{Lemma}}{}
\makeatother
\title[Tribonacci-Lucas numbers that are palindromic concatenations]{Tribonacci-Lucas numbers that are palindromic concatenations of two distinct repdigits}
\author{Mahadi Ddamulira}
\date{Source: arXiv:2509.05984v1 (CC BY 4.0)}
\begin{document}
\maketitle

\noindent\emph{Source and licence.} Tribonacci-Lucas numbers that are palindromic concatenations of two distinct repdigits. Version arXiv:2509.05984v1 (CC BY 4.0), distributed under the Creative Commons Attribution 4.0 International licence, as recorded in the arXiv licence metadata for this exact version and shown on the abstract page https://arxiv.org/abs/2509.05984v1. Full rights notice: licenses/arxiv-2509.05984-CC-BY-4.0.txt.\par\smallskip

\noindent\textit{Editorial scope.} Counted unit: the Lemma whose source label is \texttt{l4} in the article (source0.tex lines 233--240, source label \texttt{l4}), delivered together with the article's own complete original proof of it (source0.tex lines 242--272, one \texttt{proof} environment of 31 lines). No mathematics is written, repaired or re-derived: statement and proof are the article's own text line for line. Required context retained uncounted: Trib4 (lines 76--81); lem1 (lines 94--100); eq3 (lines 225--229). Every definition, display and label of those lines is retained unchanged. Omitted from this package and not redistributed: 1--75, 82--93, 101--224, 230--232, 241--241, 273--788. Nothing inside a retained excerpt is omitted or altered: every retained body is delivered exactly as the article's lines are written, so no omission receipt of any kind accompanies this package. Citations: the retained excerpts carry no citation command at all, so no bibliography entry of the article is transcribed and no reference is redistributed. Reference adaptation: none was needed and none was made; every reference inside the delivered unit resolves inside it, so no unresolved reference or citation is produced. Printed slips kept literal: no printed word, symbol, hypothesis or number of the counted statement or of its proof was altered. Layout and class: the article's own class is \texttt{article} with packages including inputenc, amsmath, url, hyperref, amsthm, bm, amsfonts, amssymb, graphicx, geometry, xcolor; the wrapper is the lane's pinned \texttt{amsart} editorial wrapper at 10pt on A4, which restates the theorem-like environments the retained text needs and adds the article's own macro definitions that the retained text needs (none), with extra wrapper packages (bm, bbm, dsfont, xspace, cancel, mathtools). The delivered numbering is the unit's own. Assets: no retained span contains a figure environment, a graphics-inclusion command, a PDF-page inclusion, a table or a TikZ picture; no image file is copied into this package, the package declares no asset dependency and no image file is redistributed with it. Licence: version arXiv:2509.05984v1 of this article is distributed under the Creative Commons Attribution 4.0 International licence, as recorded in the arXiv licence metadata for this exact version and shown on the abstract page https://arxiv.org/abs/2509.05984v1. A full rights notice is delivered at licenses/arxiv-2509.05984-CC-BY-4.0.txt. Characters: every retained excerpt is pure ASCII, so the delivered engine, XeLaTeX with its default Latin Modern text font, reports no missing character.
\begin{equation}\label{Trib4}
\begin{aligned}
1.83&<\alpha<1.84;\\
0.73 < |\beta|&=|\gamma|=\alpha^{-\frac{1}{2}}< 0.74.
\end{aligned}
\end{equation}

\begin{lemma}\label{lem1}
Let $m \ge 1$ be a positive integer. Then

\begin{equation*}
\alpha^{m-1} \leq S_m < \alpha^{m+1}.
\end{equation*}
\end{lemma}

\begin{align}\label{eq3}
S_n & = \overline{\underbrace{d_1\cdots d_1}_\text{$\ell$ times} \underbrace{d_2\cdots d_2}_\text{$m$ times} \underbrace{d_1\cdots d_1}_\text{$\ell$ times}} \nonumber \\
& = \overline{\underbrace{d_1\cdots d_1}_\text{$\ell$ times}} \times 10^{\ell + m} + \overline{ \underbrace{d_2\cdots d_2}_\text{$m$ times}} \times 10^\ell + \overline{\underbrace{d_1\cdots d_1}_\text{$\ell$ times}} \nonumber \\
\Rightarrow ~S_n & = \frac{1}{9} \left(d_1 \times 10^{2\ell + m} - (d_1 - d_2) \times 10^{\ell+m} + (d_1-d_2) \times 10^\ell -  d_1 \right).
\end{align}

\begin{lemma}\label{l4}
All solutions of the Diophantine equation \eqref{eq3} satisfy the inequality below

\begin{equation*}
(2\ell + m) \log 10 - 3 < n \log \alpha < (2\ell + m) \log 10 + 1.
\end{equation*}

\end{lemma}

\begin{proof}
\noindent From Lemma  \ref{lem1}, we have $\alpha^{n-1} \le  S_n <  \alpha^{n+1}$. So, we note that 
\begin{equation*}
\alpha^{n-1} \le  S_n < 10^{2\ell + m}.
\end{equation*}

\noindent Taking the logarithm on both sides of the preceeding inequality and simplifying,  we get 

\begin{equation*}
n \log \alpha < (2\ell + m) \log 10 + \log \alpha.
\end{equation*}

\noindent Hence, using the estimates in \eqref{Trib4}, we get
\begin{equation}\label{qw1}
n \log \alpha < (2\ell + m) \log 10 + 1. 
\end{equation}
\noindent
On the other hand, using Lemma \ref{lem1},  we also note that 
\begin{align*}
10^{2\ell + m -1} < S_n <  \alpha^{n+1}.
\end{align*}
Taking logarithms on both sides and simplifying as before gives
\begin{align*}
(2\ell +m)\log 10-\log 10 - \log \alpha < n\log \alpha.
\end{align*}
Similary, using the estimates in \eqref{Trib4}, we get that
\begin{equation}\label{qw2}
(2\ell +m)\log 10-3 < n\log \alpha.
\end{equation}
Clearly, from \eqref{qw1} and \eqref{qw2}, we get the desired inequality in Lemma \ref{l4}, as required.
\end{proof}

\end{document}
